\section*{Hoofdstuk \ref{ch:b2:phanerozoic-earth} --- De fanerozoïsche Aarde}

\begin{solution}{exo:b2:phanerozoic-earth:1}
(a) Parijs (\ang{5.1}~NB) en New York (\ang{3.6}~ZB); Kaapstad en Sydney lagen dicht bij \ang{60}~ZB.
(b) Kaapstad: $61.9 - 33.9 = \ang{28.0}$, $28.0 \times 111.2 = \qty{3114}{km}$; New York:
$40.7 + 3.6 = \ang{44.3}$, \qty{4926}{km}. (c) $3114/300 = \qty{10.4}{km}$ per miljoen jaar,
\qty{1.0}{cm} per jaar; $4926/300 = \qty{16.4}{km}$ per miljoen jaar, \qty{1.6}{cm} per jaar.
\end{solution}

\begin{solution}{exo:b2:phanerozoic-earth:2}
(a) 72, 120.1, 50.5 en \qty{135.4}{Myr}; gemiddelde $(444 - 66)/4 = \qty{94.5}{Myr}$. (b) De
crisis aan het einde van het Perm (Siberië, \qty{1.9}{Myr} ertussen), die aan het einde van het Trias (CAMP, \qty{0.4}{Myr}) en die aan het
einde van het Krijt (Deccan, \qty{0.5}{Myr}). (c) Nee: de intervallen lopen van 50 tot
\qty{135}{Myr}; een gemiddelde bestaat voor elke reeks data.
\end{solution}

\begin{solution}{exo:b2:phanerozoic-earth:3}
(a) $\qty{91}{m}/\qty{0.036}{mm.yr^{-1}} = \num{2.5e6}$ jaar. (b) $3.7/0.036 = 103$ keer. (c)
Ver van elke kust, in helder water, bereikte weinig modder uit rivieren de bodem; de regen van
algenschaaltjes was het voornaamste sediment.
\end{solution}

\begin{solution}{exo:b2:phanerozoic-earth:4}
(a) Van het rivierzandsteen tot het krijt worden de afzettingen naar boven dieper: een transgressie. Het
strandzand boven het krijt legt een regressie vast. (b) Nee: ze legt een stijging van het relatieve zeeniveau vast,
die ook door daling van de plaats, of door een afname van de sedimentaanvoer, kan worden veroorzaakt. (c)
Dezelfde transgressie, op hetzelfde tijdstip, op andere continenten vinden, of de
daling verwijderen door \omterm{def:b2:basins-compared:backstripping}{backstripping}.
\end{solution}

\begin{solution}{exo:b2:phanerozoic-earth:5}
(a) $233/97 = \qty{2.4}{m}$ per miljoen jaar, \qty{2.4}{\micro m} per jaar. (b) $3.7/0.0024 =
1540$: de huidige stijging is zo'n 1500 keer sneller. (c) De gebackstripte curve: boringen hebben aangetoond
dat cyclusgrafieken de timing van dalingen goed weergeven maar hun grootte overschatten, met een factor twee of
meer; New Jersey heeft zijn eigen vertekening (dynamische topografie), vandaar de schatting van $100 \pm
\qty{50}{m}$.
\end{solution}

\begin{solution}{exo:b2:phanerozoic-earth:6}
(a) $\num{e5} \times \num{e9} \times 2900 = \qty{2.9e17}{kg}$ basalt, en $\num{2.9e17}
\times 0.0075 = \qty{2.2e15}{kg}$ koolstofdioxide, $\num{2.2e15}/0.04401 = \qty{4.9e16}{mol}$.
(b) $\num{278e-6} \times (44.01/28.97) \times \num{5.1e18} = \qty{2.15e15}{kg}$. (c) Eén puls
gaf ongeveer evenveel af als de hele lucht in 1750 bevatte. (d) $\num{e17}/\num{2.15e15} = 46.5$: bijna vijftig.
\end{solution}

\begin{solution}{exo:b2:phanerozoic-earth:7}
(a) $2630 \times \frac{\pi}{6} \times (\num{e4})^3 = \qty{1.38e15}{kg}$. (b) $\tfrac12 \times
\num{1.38e15} \times (\num{11.2e3})^2 = \qty{8.7e22}{J}$; bij \qty{20}{km.s^{-1}},
\qty{2.76e23}{J}. (c) Een jaar van \qty{47}{TW} is $\num{47e12} \times \num{3.156e7} =
\qty{1.48e21}{J}$: 59 en 186 jaar.
\end{solution}

\begin{solution}{exo:b2:phanerozoic-earth:8}
(a) Donkere kleur, fijne laminatie, geen graafgangen, veel organisch materiaal, pyriet; vaak een verschuiving van de
koolstofisotopen en tekenen van opwarming. (b) De toarcische gebeurtenis met Karoo--Ferrar, OAE~1a met
Ontong Java. (c) Data hebben onzekerheden tot miljoenen jaren, en er bestaan veel gebeurtenissen en
provincies, dus toevallige samenvallingen zijn waarschijnlijk; een oorzakelijk verband vraagt dat de provincie
net vóór de gebeurtenis begint en een tracer van het vulkanisme, zoals kwik, in dezelfde sedimenten.
\end{solution}

\begin{solution}{exo:b2:phanerozoic-earth:9}
(a) $60.9 + 5.1 = \ang{66.0}$, $66.0 \times 111.2 = \qty{7339}{km}$ in \qty{200}{Myr}:
\qty{36.7}{km} per miljoen jaar, \qty{3.7}{cm} per jaar. (b) Meer dan twee keer de
\qty{1.6}{cm} per jaar sinds \qty{300}{Ma}. (c) Alleen de verandering van breedte is bekend; de verandering van
lengte, vrij in een paleomagnetisch referentiekader, en elke omweg vergroten de afstand.
\end{solution}

\begin{solution}{exo:b2:phanerozoic-earth:10}
(a) De fractie bodem met ouderdom $t$ is $p(t) = \frac2T(1 - t/T)$. Dan is $\int_0^T tp\,\mathrm{d}t =
\frac2T\bigl(\frac{T^2}2 - \frac{T^2}3\bigr) = \frac T3$ en $\int_0^T \sqrt t\,p\,\mathrm{d}t =
\frac2T\bigl(\frac23 T^{3/2} - \frac25 T^{3/2}\bigr) = \frac8{15}\sqrt T$. (b) $T = 3 \times 64.2 =
\qty{192.6}{Myr}$; $2500 + \frac8{15} \times 350 \times 13.878 = \qty{5091}{m}$. (c) $T/2 =
\qty{96.3}{Myr}$: \qty{4332}{m}; de daling is \qty{759}{m}, \qty{2}{\percent} minder dan
\qty{774}{m}. (d) Gelijkmatig model: $\frac23\sqrt\tau = \frac23\sqrt{2\bar t} = 0.943\sqrt{\bar t}$;
driehoek: $\frac8{15}\sqrt{3\bar t} = 0.924\sqrt{\bar t}$. Het antwoord hangt veel meer van de gemiddelde ouderdom
af dan van de vorm.
\end{solution}

\begin{solution}{exo:b2:phanerozoic-earth:11}
(a) In $\Delta s$ hangt alleen de factor $1 - f^{-1/2}$ af van $f$; deel door zijn waarde bij $f = 2$,
$1 - 2^{-1/2} = 0.2929$. (b) $f = (1 - 0.2929\,\Delta s/432)^{-2}$: 1.071 voor \qty{50}{m}, 1.239
voor \qty{150}{m}; 7 tot \qty{24}{\percent} sneller dan vandaag. (c) $432 \times (1 -
0.8^{-1/2})/0.2929 = \qty{-174}{m}$: de zee zou \qty{174}{m} dalen. (d) De ouderdomsverdeling
verandert alleen naarmate oude bodem wordt gesubduceerd en vervangen, wat een tijd van de orde van $\tau$ vergt,
ongeveer \qty{100}{Myr}; de zee volgt over tientallen miljoenen jaren.
\end{solution}

\begin{solution}{exo:b2:phanerozoic-earth:12}
(a) $2^{-7.4/2} = 0.077$: \qty{92}{\percent} verdwenen. (b) $t = 2\log_2(1/0.19) = \qty{4.8}{Myr}$.
(c) $4.8/0.060 = 80$ keer sneller. (d) Binnen een lange etage verwijdert achtergronduitsterving alleen al
de meeste soorten, dus een korte puls telt op bij een groot totaal en is er moeilijk van te scheiden; de
onvolledigheid van het archief smeert de laatste verschijningen bovendien naar achteren uit.
\end{solution}

\begin{solution}{pb:b2:phanerozoic-earth:1}
\textbf{1.} 2500, 4250, 5304 en \qty{6000}{m}.
\textbf{2.} $\tau = 2 \times 64.2 = \qty{128.4}{Myr}$; $2500 + 350\sqrt{128.4} = \qty{6466}{m}$.
\textbf{3.} $\num{293e6}/(65000 \times 128.4) = \qty{35.1}{km}$ per miljoen jaar,
\qty{35.1}{mm} per jaar: een trage snelheid.
\textbf{4.} $\frac1\tau\int_0^\tau(2500 + 350\sqrt t)\,\mathrm{d}t = 2500 + \frac23 \times
350\sqrt\tau = \qty{5144}{m}$.
\textbf{5.} De wortelwet is te diep voor bodem ouder dan enkele tientallen miljoenen jaren,
die afvlakt; de echte ouderdommen zijn niet gelijkmatig verdeeld, omdat jonge bodem het meest voorkomt; en
sediment, \omterm{def:b2:oceanic-islands:seamount}{zeebergen} en plateaus verhogen de bodem.
\textbf{6.} $\num{293e6} \times 5.144 = \qty{1.51e9}{km^3}$, \qty{13}{\percent} meer dan de
hele oceaan: het model is te diep.
\textbf{7.} $v = \qty{17.6}{km}$ per miljoen jaar; $17.6 \times 128.4 = \qty{2260}{km}$ aan
elke kant.
\textbf{8.} $t = 1000/17.6 = \qty{56.8}{Myr}$; $2500 + 350\sqrt{56.8} = \qty{5138}{m}$.
\textbf{9.} Op één flank, met $x = vt$: $\int_0^{v\tau}\bigl(d(\tau) - d(x/v)\bigr)\mathrm{d}x
= vc\int_0^\tau(\sqrt\tau - \sqrt t)\,\mathrm{d}t = \frac13 vc\,\tau^{3/2}$; het dubbele voor beide.
$\frac23 \times 17.6 \times 0.350 \times 128.4^{3/2} = \frac23 \times 17.6 \times 0.350 \times
1454.9 = \qty{5975}{km^2}$.
\textbf{10.} $65000 \times 5975 = \qty{3.88e8}{km^3}$.
\textbf{11.} $\num{3.88e8}/\num{1.33e9} = 0.29$.
\textbf{12.} Zijn bodem is gemiddeld jonger, dus heter, minder dicht en hoger; de band van jonge
bodem langs de as is breder.
\textbf{13.} $\tau = \qty{64.2}{Myr}$; $2500 + \frac23 \times 350\sqrt{64.2} = \qty{4370}{m}$.
\textbf{14.} $5144 - 4370 = \qty{774}{m}$.
\textbf{15.} $\num{293e6} \times 0.774 = \qty{2.27e8}{km^3}$.
\textbf{16.} $\num{2.27e8}/\num{362e6} = \qty{0.627}{km}$, \qty{627}{m}.
\textbf{17.} $\Delta s = 627 \times (3300 - 1025)/3300 = \qty{432}{m}$.
\textbf{18.} $f = \bigl(1 - 0.2929 \times 100/432\bigr)^{-2} = 1.15$: \qty{15}{\percent} sneller.
\textbf{19.} $0.149 \times 148 = \qty{22.1}{\mega km^2}$.
\textbf{20.} $48.1 + 0.32 \times 7.7 = \qty{50.6}{\percent}$, \qty{74.9}{\mega km^2}.
\textbf{21.} Terwijl ze stijgt, bedekt ze meer oppervlakte, dus hetzelfde volume vormt een dunnere laag
(\cref{prop:b2:phanerozoic-earth:flooding}).
\textbf{22.} Overstroomde oppervlakten bij 100, 200, 300 en \qty{400}{m}: 22.1, 41.7, 57.3 en
\qty{71.2}{\mega km^2}. Stappen: $(362 + 11.0) \times 0.1 = \qty{37.3e6}{km^3}$, daarna 39.4, 41.2
en \qty{42.6e6}{km^3}; gecumuleerd 37.3, 76.7, 117.9 en \qty{160.5e6}{km^3}. Het te plaatsen volume
is \qty{156.4e6}{km^3}, tussen 300 en \qty{400}{m}: $300 + (156.4 - 117.9)/42.6 \times 100
= \qty{390}{m}$.
\textbf{23.} $38.7 + 0.90 \times 9.4 = \qty{47.2}{\percent}$ van het huidige land.
\textbf{24.} Een verdubbeling van de spreidingssnelheid zou de zee ongeveer \textbf{\qty{390}{m}} doen stijgen en
bijna de helft van het huidige land laten verdrinken.
\end{solution}
